\chapter*{ภาคผนวก ข: เอกลักษณ์แคลคูลัสเวกเตอร์และตัวดำเนินการ}
\addcontentsline{toc}{chapter}{ภาคผนวก ข: เอกลักษณ์แคลคูลัสเวกเตอร์และตัวดำเนินการ}

รวบรวมเอกลักษณ์สำคัญของพีชคณิตและแคลคูลัสเวกเตอร์ เพื่อประโยชน์ในการวิเคราะห์สมการเชิงอนุพันธ์ย่อยในทฤษฎีสนามแม่เหล็กไฟฟ้าและกลศาสตร์ควอนตัม

\vspace{0.4cm}
\section*{1. เอกลักษณ์เวกเตอร์พื้นฐาน}
\begin{align}
\mathbf{A} \cdot (\mathbf{B} \times \mathbf{C}) &= \mathbf{B} \cdot (\mathbf{C} \times \mathbf{A}) = \mathbf{C} \cdot (\mathbf{A} \times \mathbf{B}) \\
\mathbf{A} \times (\mathbf{B} \times \mathbf{C}) &= (\mathbf{A} \cdot \mathbf{C})\mathbf{B} - (\mathbf{A} \cdot \mathbf{B})\mathbf{C}
\end{align}

\vspace{0.2cm}
\section*{2. อนุพันธ์อันดับหนึ่งของผลคูณ}
\begin{align}
\nabla(fg) &= f\nabla g + g\nabla f \\
\nabla(\mathbf{A} \cdot \mathbf{B}) &= \mathbf{A} \times (\nabla \times \mathbf{B}) + \mathbf{B} \times (\nabla \times \mathbf{A}) + (\mathbf{A} \cdot \nabla)\mathbf{B} + (\mathbf{B} \cdot \nabla)\mathbf{A} \\
\nabla \cdot (f\mathbf{A}) &= f(\nabla \cdot \mathbf{A}) + \mathbf{A} \cdot \nabla f \\
\nabla \cdot (\mathbf{A} \times \mathbf{B}) &= \mathbf{B} \cdot (\nabla \times \mathbf{A}) - \mathbf{A} \cdot (\nabla \times \mathbf{B}) \\
\nabla \times (f\mathbf{A}) &= f(\nabla \times \mathbf{A}) + (\nabla f) \times \mathbf{A} \\
\nabla \times (\mathbf{A} \times \mathbf{B}) &= (\mathbf{B} \cdot \nabla)\mathbf{A} - (\mathbf{A} \cdot \nabla)\mathbf{B} + \mathbf{A}(\nabla \cdot \mathbf{B}) - \mathbf{B}(\nabla \cdot \mathbf{A})
\end{align}

\vspace{0.2cm}
\section*{3. อนุพันธ์อันดับสองและทฤษฎีบทอินทิกรัล}
\begin{align}
\nabla \times (\nabla f) &= 0 \quad (\text{Curl of Gradient is Zero}) \\
\nabla \cdot (\nabla \times \mathbf{A}) &= 0 \quad (\text{Divergence of Curl is Zero}) \\
\nabla \times (\nabla \times \mathbf{A}) &= \nabla(\nabla \cdot \mathbf{A}) - \nabla^2 \mathbf{A}
\end{align}

\vspace{0.2cm}
\noindent\textbf{ทฤษฎีบทไดเวอร์เจนซ์ (Divergence Theorem):}
\begin{equation}
\iiint_V (\nabla \cdot \mathbf{A}) \, dV = \iint_S \mathbf{A} \cdot d\mathbf{S}
\end{equation}

\vspace{0.2cm}
\noindent\textbf{ทฤษฎีบทของสโตกส์ (Stokes' Theorem):}
\begin{equation}
\iint_S (\nabla \times \mathbf{A}) \cdot d\mathbf{S} = \oint_C \mathbf{A} \cdot d\mathbf{r}
\end{equation}
